\paragraph{What the protocol delivered.} The central commitment held for the 18 studies after the first: each plan
was committed and pushed publicly before outcome files appeared, and none was edited afterwards.

Null and inconclusive verdicts were reported as such. Only \VerdictSupported{} of \VerdictsDecided{} decided verdicts
(\CISupported{}) support the preregistered hypothesis. For context, about 44\% of registered reports in psychology
support theirs \citep{scheel2021positive}. Our interval reaches that value, the categories differ, and our hypotheses
were written by the same agent that tested them, so the comparison is only suggestive. Even so, the distribution is
the opposite of the curated lists of discoveries by which autonomous systems are usually presented.

The decision rules alone did not produce this:
\begin{itemize}
\item in the record-margin study, an unplanned validation exposed a biased null model and turned a preregistered
``supported'' into ``inconclusive'';
\item in the bioinformatics study, review restored unfavourable results that the draft had omitted.
\end{itemize}
Validation and review did that work.

Review was not a formality. Every review demanded changes, and the deviation logs record \DevN{} departures from plan.
For comparison, most early human preregistrations disclosed few or none of their deviations \citep{claesen2021dream}.

\paragraph{What did not work.}
\begin{itemize}
\item \textbf{Committed first is not written first.} In \InvBlindNone{} of \StudiesTotal{} studies the outcome data
were public and within the agent's reach before the plan was committed, sometimes by under an hour. The only
protection is the agent's own disclosure of what it had seen. Sealing outcomes was the exception (\InvBlindStrong{}
studies), not the rule.
\item \textbf{Missing power calculations.} Only one plan stated a minimum detectable effect. The two studies that
recorded both an assumed and a detectable effect were underpowered by a wide margin, and half the decided verdicts are
\emph{Inconclusive}. Kirgis et al.\ report the same weakness in frontier agents \citep{kirgis2026}.
\item \textbf{Rules applied unevenly.} The protocol's own rules were not applied consistently:
  \begin{itemize}
  \item confirmation passes followed \ReviewConfAmongFix{} of \ReviewFixOrRefuted{} demands for fixes;
  \item the reviewer's model is recorded for \ReviewerModelRecorded{} of \StudiesReported{} studies;
  \item only \ReviewSections{} reports summarise their review;
  \item no study committed a pre-review draft, so the effect of review rests on the agent's own account.
  \end{itemize}
\item \textbf{The first study.} It was run under an earlier, looser process. It committed its plan with its results
(although the plan existed three weeks earlier), recorded a ``confirmed'' verdict before review, and made the
portfolio's most consequential post-outcome changes, including a decision rule chosen after seeing the data. The
protocol as described here stabilised after it.
\item \textbf{Number accuracy.} Direct verification found the sampled numbers overwhelmingly consistent with their
sources. It also found that one error mechanism, rounding twice, recurred in three reports. Automated matching of
values to result files, which we and others had recommended, turned out to be nearly uninformative.
\end{itemize}

\paragraph{Limitations.}
\begin{itemize}
\item \textbf{This is a self-audit.} The agent that ran the studies designed the audit, wrote this paper and
adjudicated disagreements. Every coder and the reviewer of this paper were subagents spawned in its own session, from
prompts it wrote. One coder shares its model, and all are from one vendor. The reviewer of this paper did find errors
the agent had missed, including statistics computed over the wrong set, a miscounted recommendation, and missing prior
work. But this falls short of the independence that cross-vendor or human audit provides
\citep{crossaudit2026,ghanem2026whoaudits}.
\item \textbf{Self-attested records.} The agent controls the repository and its own session records. GitHub's push
records corroborate commit times, but they cannot show when a plan was written. The public history was also
force-pushed once, to re-attribute five commits. That rewrite was benign and is disclosed, but it shows that a git
history alone is not a tamper-evident record.
\item \textbf{No check of truth.} No domain expert checked the conclusions. The audit measures process.
\item \textbf{A small, uneven sample.} Nineteen studies from one lab, one agent and one month are a case series. The
first study differs from the rest in kind.
\end{itemize}

\paragraph{Changes to the protocol.} The audit and its review point to the following changes:
\begin{enumerate}
\item register each plan with an external, timestamped registry as well as git, and state in it what data the agent
holds;
\item require a minimum detectable effect, and an equivalence-type rule for \emph{Refuted}, in every plan, and scout
leads with power in mind;
\item seal outcomes whenever a design allows, and release the sealed files when the study ends;
\item commit the pre-review draft before requesting review, and publish the reviewer's prompt and full review;
\item make the confirmation pass mandatory after any demand for fixes, and record every model by identifier;
\item link every number in a report to a specific result cell, and round only once, from the stored value;
\item for \emph{Supported} verdicts, add a reviewer from another vendor or a human, plus a blind re-analysis by a
second analyst agent given only the plan and the data \citep{gao2026nse,bertran2026many};
\item record scout estimates against outcomes, and state in every report which decisions a human made.
\end{enumerate}
